% !TEX root = ../../main.tex
% C4.1 - Actors (translated). First place that names the roles: Administrator, Operator, Viewer; Case.
\section{Actors}
\label{sec:4.1}

The system has three actors, corresponding to three kinds of user account divided by responsibility: the \textbf{Administrator} (Admin), the \textbf{Operator} and the \textbf{Viewer}. All three share one login function; after authentication, the system determines the role of the account and only offers the functions and data suited to that role.

Before describing each actor, two business concepts used throughout the following chapters need to be defined:
\begin{itemize}
    \item \textbf{Surveillance area:} a geographic scope such as a building or a public space. Each camera belongs to exactly one area, set when the camera is added and not changed afterwards. The list of areas is declared at deployment; the current version has no interface for adding, editing or deleting areas.
    \item \textbf{Case:} a record created by an Operator to keep the search results related to an incident, with a title, notes and a status of \emph{Open} or \emph{Closed}. A case belongs to the Operator in charge of it but is not tied directly to a surveillance area: the area only limits the scope of search, while the rights over a case follow its owner. Since an Operator's area can change, a case may contain results from several areas.
\end{itemize}

\subsection{Administrator}
\label{sec:4.1.1}

The Administrator is responsible for the \emph{technical administration} of the system, keeping the camera system and the AI processing running so that the other actors can use them. The Administrator's responsibilities are:
\begin{itemize}
    \item \textbf{Managing accounts} of Operators and Viewers, including assigning a surveillance area to each Operator.
    \item \textbf{Managing cameras} and the connection to the RTSP stream of each camera.
    \item \textbf{Managing AI processing:} deciding which cameras are analysed, and choosing the person detector (Detector) and the tracking algorithm (Tracker) used by the whole system.
    \item \textbf{Monitoring operation:} following the status of cameras and AI processing, checking the AI components, and consulting the audit log.
\end{itemize}

Technical administration rights do \emph{not} imply business rights: by default, the Administrator may not search for people and does not see cases. The Administrator account is created at deployment; through the interface, the Administrator can only create and assign the Operator and Viewer roles.

\subsection{Operator}
\label{sec:4.1.2}

The Operator uses the main business function of the system directly: searching for people in the analysed camera data and saving the matching results into cases. The Operator's responsibilities are:
\begin{itemize}
    \item \textbf{Searching for people} by image, English description or appearance attributes.
    \item \textbf{Assessing results:} viewing each result in the context of its frame and deciding which results really show the person being sought.
    \item \textbf{Managing cases:} saving matching results into cases and following each case until it is closed.
\end{itemize}

The Operator's scope is limited by two principles. First, the \emph{search scope} is set by the surveillance area attached to the account: the system only searches the operational cameras of that area, and rejects any request containing a camera or area outside the Operator's rights. Second, the \emph{case scope} is set by ownership: each case has exactly one Operator in charge, the one who created it; an Operator only acts on the cases he or she owns and cannot transfer ownership to someone else. The two scopes are independent: when moved to another area, an Operator still sees and manages his or her old cases, while new searches only run in the new area.

\subsection{Viewer}
\label{sec:4.1.3}

The Viewer is at the management level of the organisation, such as a head of department or a director, who needs to follow the results of the work but does not search or configure the system. The Viewer's responsibilities are:
\begin{itemize}
    \item \textbf{Viewing an overview} of the cases across the whole system.
    \item \textbf{Viewing cases:} every case and the results saved in it.
\end{itemize}

The Viewer only has read access and does not search for people. Unlike the Operator, the Viewer is not limited by area, because the Viewer's scope is the cases, which are not tied to areas.

\subsection{Principles of Role Separation}
\label{sec:4.1.4}

The division above follows two principles: \textbf{separating technical rights from business rights}, so that the people who configure the system do not automatically access search data and the people who use the data cannot interfere with the configuration; and \textbf{least privilege}, so that each actor only has the rights needed for its work. The rights of each role for each group of actions are summarised in the permission matrix of Section~\ref{sec:4.6}.

The camera streams and the background AI processing are not actors: a camera is a data source configured by the Administrator, and the analysis of camera data is performed automatically by the system when AI processing is enabled, not triggered directly by a user.
